\cotefile{157-4}
\scanpage{273}
\mspage{\unsure{4|5}}

\deleted{(pas de page 4)}

\deleted{\uline{Loc} $(3')$ : Soit $f:X\longrightarrow S$ un foncteur Cat-\unsure{fibrant|fibré} \unsure{de|dans} Cat, \deleted{\unsure{un}} telle que les fibres $X_s$ soient $W$-asphériques, alors $f\in W$.}

\deleted{\uline{Loc} $(3'\,\mathrm{bis})$ Soit \deleted{$f:X\unsure{\longrightarrow|{-}}Y$ \unsure{un}} \inserted{$S$ dans Cat,} $X,Y$ deux cat. Cat-fibrées scindées}

\uline{Loc} $(3')$ \textcircled{a} \inserted{(Axiome des Cat-fibrations, forme \inserted{\unsure{disons|dite}} faible)} \deleted{\unsure{Soit}}\unsure{Soient|Soit} $\deleted{\unsure{X|D}}S$ dans Cat, $X$ une catégorie fibrée scindée sur $S$, telle que les \inserted{fibres} $X_s$ $(s\in\operatorname{Ob}S)$ soient $W$-\unsure{asphériques}. Alors $X\unsure{{-}|\longrightarrow} S$ est dans $W$

\textcircled{b} Soit $X$ dans $\mathrm{Cat}_{/S}$, alors $\unsure{i_X|\iota_X}:X\longrightarrow\Phi_S(X)$ est dans $W$

% slanted
\margin{b) \unsure{concerne} la \uline{comparaison} d'une $X$ quelc. sur $S$, avec la cat. fibrée scindée $\Phi_S(X)$ sur $S$}

\uline{Loc} $(3'\,\mathrm{bis})$ : (Axiome des Cat-fibrations, forme \unsure{disons|dite} forte) : \textcircled{a} \unsure{Soient|Soit} $S$ dans Cat, $X,Y$ deux catégories fibrées scindées sur $S$, $f:X\longrightarrow Y$ un $S$-foncteur cartésien compatible aux scindages, tel que $\forall s\in\operatorname{Ob}S$, le foncteur induit sur les fibres

\[
f_s:X_s{-}Y_s
\]

soit dans $W$. Alors $f\in W$.

\textcircled{b} Comme \deleted{Loc} $(3')$ b).

\uline{NB} En fait, moyennant $\operatorname{Loc}(1,2a)$, les énoncés $\operatorname{Loc}(3')$, $\operatorname{Loc}(3'\,\mathrm{bis})$ impliquent les énoncés similaires sans scindages, et sans même supposer $f$ cartésien dans $\operatorname{Loc}(3'\,\mathrm{bis})$, cf démonstration ci-dessous.

\uline{Démonstration} \qquad a) $\operatorname{Loc}(1,2,3)\Longrightarrow\operatorname{Loc}(1,2a,3')$

Il suffit de prouver $\operatorname{Loc}(3')$. \inserted{Prouvons $\operatorname{Loc}(3'\,a)$.} Pour ceci, sous les conditions de $\operatorname{Loc}(3'\,a)$, prouvons que $X\longrightarrow S$ est $W$-\inserted{co}asphérique, i.e. que les catégories

\[
{}_{s\backslash}X=X\times_S{}_{s\backslash}S
\]

sont $W$-asphériques. Il en résultera bien que $f\in W$ en vertu de $\operatorname{Loc}(3)$ \inserted{appliqué à $f^{o}$,} et de $W=W^{o}$. Mais

% archive sheet number: 272

\scanpage{274}
\mspage{P 30}

\deleted{``discrètes'' associées aux différents nombres premiers, devaient fournir la clef pour un renouvellement de vaste envergure de l'arithmétique. C'est dans cet esprit qu'il a dégagé, en 1949, les célèbres ``c o n j e c t u r e s \quad d e \quad W e i l''. Conjectures absolument époustouflantes, à vrai dire, qui faisaient entrevoir, pour ces nouvelles ``variétés'' (ou ``espaces'') de nature discrète, la possibilité de certains types de constructions et d'arguments (*)\unsure{,|\mbox{}} qui jusque là ne semblaient pensables que dans le cadre des seuls ``espaces'' considérés comme dignes de ce nom par les analystes - savoir, les espaces dits ``topologiques'' (où la notion de variation continue a cours).}

\deleted{On peut considérer que la géométrie nouvelle est,avant toute autre chose, une s y n t h è s e entre ces deux mondes, jusque là mitoyens et étroitement solidaires, mais pourtant séparés : le m o n d e \quad ``a r i t h m é t i q u e'', dans lequel vivent les (soi-disants) ``espaces'' sans principe de continuité, et le m o n d e \quad d e \quad l a \quad g r a n d e u r \quad c o n t i n u e, ou vivent les ``espaces'' au sens propre du terme, accessibles aux moyens de l'analyste et (pour cette raison même) acceptés par lui comme dignes de giter dans la cité mathématique. D a n s \quad l a \quad v i s i o n \quad n o u v e l l e, \quad c e s \quad d e u x \quad m o n d e s \quad j a d i s \quad s é p a r é s, \quad n ' e n \quad f o r m e n t \quad p l u s \quad q u ' u n \quad s e u l.}

\deleted{Le premier embryon de cette vision d'une ``g é o m é t r i e \quad a r i t h m é t i q u e'' (comme je propose d'appeler cette géométrie nouvelle) se trouve dans les conjectures de Weil. Dans le développement de certains de mes thèmes principaux (**) , ces conjectures sont restées ma principale source d'inspiration, tout au long des années entre 1958 et 1969. Dès avant moi, d'ailleurs, O s c a r \quad Z a r i s k i d'un côté, puis J e a n - P i e r r e \quad S e r r e de l'autre, avaient développé pour les espaces-sans-foi-ni-loi de la géométrie algébrique ``abstraite'' certaines méthodes ``topologiques'', inspirées de celles qui avaient cours précédemment pour les ``espaces bon teint'' de tout le monde (***).}

\rule{2cm}{0.4pt}

\deleted{(*) (A l'intention du lecteur mathématicien.) Il s'agit ici\unsure{.|\mbox{}} des ``constructions et arguments'' liés à la théorie cohomologique des variétés différentiables ou complexes, et notamment de ceux impliquant la formule des points fixes de Lefschetz, et la théorie de Hodge.}

\deleted{(**) Il s'agit des quatre thèmes ``médians'' (n\textsuperscript{o}s 5 à 8), savoir ceux des t o p o s de la c o h o m o l o g i e \quad é t a l e et \handwritten{$\ell$}-adique, des m o t i f s, et (dans une moindre mesure)\unsure{'|\mbox{}} celui des c r i s t a u x. J'ai dégagé ces thèmes tour à tour entre 1958 et 1966.}

\deleted{(***) (A l'intention du lecteur mathématicien.) La principale contribution de Zariski dans ce sens me paraît l'introduction de la ``topologie de Zariski'' (qui plus tard a été un outil essentiel pour Serre dans FAC), et son ``principe de connexité'' et ce qu'il a appelé sa ``théorie des fonctions holomorphes'' - devenus entre \deleted{\unsure{s|m}}mes mains la théorie des schémas formels, et les ``théorèmes de comparaison''}

% archive sheet number: 273

\scanpage{275}
\mspage{5}

\deleted{Mais on a}

\deleted{On est ramené à prouver ceci.}

\deleted{\uline{Lemme} 1.4. Supposons $\operatorname{Loc}(1,2,3)$, et soit $X$ fibrée sur $S$, aux fibres $W$-asphériques, $S$ ayant un objet initial. Alors $X$}

considérons l'inclusion canonique

\[
i_s:X_s\hookrightarrow{}_{s\backslash}X,
\]

il résulte du fait que $p:X\unsure{{-}|\longrightarrow}S$ est \inserted{Cat-}fibrant que ce foncteur est $W$-asphérique, (il a même un adjoint à droite i.e. les catégories $X_s/u$, où $u\in\operatorname{Ob}({}_{s\backslash}X)$, ont un objet final).

Donc par $\operatorname{Loc}(3)$ \deleted{il est} on a $i_s\in W$, et comme $X_s$ est $W$-asph par hyp., de $\unsure{=}\;{}_{s\backslash}X$.

\inserted{Prouvons $\operatorname{Loc}(3'\!\unsure{!|{}}\,b)$\inserted{\unsure{cf.|cit.}} : on sait que \inserted{$\unsure{i_X|\iota_X}$} $X\longrightarrow\Phi_S(X)$ est \inserted{$W$-}\uline{coasphérique} (il a $=$ un adjoint à gauche), donc $i_X^{o}$ est $W$-asphérique donc $\in W$, donc $i_X\in W$ car $W=W^{o}$.}

\[
\text{a bis)}\qquad \operatorname{Loc}(1,\unsure{2|1},3\,\mathrm{bis})\Longrightarrow \operatorname{Loc}(1,2,3'\,\mathrm{bis})
\]

Cette fois il faut prouver \inserted{seulement} $\operatorname{Loc}(3'\,\mathrm{bis})$. \inserted{Pour b) \deleted{\unsure{on l'a}} c'est déjà fait.} \deleted{\unsure{Soit|Si}}\unsure{Reste} \inserted{à prouver a).} $f$ est comme dans l'énoncé \inserted{$\operatorname{Loc}(3'\,\mathrm{bis}\ a)$}, prouvons que pour \deleted{\unsure{tout}} \inserted{\unsure{des|les}} $s\in S$, le \deleted{\unsure{fonct}} morphisme induit

\[
{}_{s\backslash}f:{}_{s\backslash}X\unsure{\longrightarrow|\longleftrightarrow|\hookrightarrow}{}_{s\backslash}Y
\]

est dans $W$. Comme $W=W^{o}$, on aura \unsure{donc|des} $({}_{s\backslash}f)^o\in W$, et appliquant $\operatorname{Loc}(3,\mathrm{bis})$ à $f^o$, on \unsure{trouve|aura} $f^o\in W$, donc $f\in W$, et on aura fini. Mais $ {}_{s\backslash}f$ s'insère dans \deleted{$\unsure{{}_{s\backslash}X\longleftarrow}$} le carré commutatif

% archive sheet number: 274

\scanpage{276}
\mspage{6}

\diagram{\begin{tikzcd}[ampersand replacement=\&]
X_s \arrow[r,dash,"f_s","{\unsure{{-}|\longrightarrow}}"{pos=0.97,description}] \arrow[d,dash,"i_s^X"'] \& Y_s \arrow[d,dash,densely dashed,"i_s^Y"]\\
{}_{s\backslash}X \arrow[r,dash,"{{}_{s\backslash}f}"] \& {}_{s\backslash}Y,
\end{tikzcd}}

où, on l'a vu \unsure{dans|lors} a), $i_s^X$ et $i_s^Y$ sont dans $W$. Par hypothèse on a $f_s\in W$, donc $ {}_{s\backslash}f\in W$, \unsure{qed}.

\[
\text{b)}\qquad \operatorname{Loc}(1,2a,3')\Longrightarrow\operatorname{Loc}(1,2,3)
\]

Prouvons d'abord $\operatorname{Loc}(2)$. Appliquant $\operatorname{Loc}(3'\,a)$ à la projection

\[
\unsure{\mathbb{\Delta}|\Delta}^{1}\times X{-}X,\qquad(\text{où }X\in\operatorname{Ob}\mathrm{Cat}),
\]

\deleted{\unsure{et}} qui est fibrée scindée à fibres isom. à $\mathbb{\Delta}^{1}$, \deleted{\unsure{dont}} $W$-asphériques par $\operatorname{Loc}(2a)$, on voit que cette projection est dans $W$. Cela permet d'appliquer les arguments \unsure{standart|standard} d'homotopie (cf $\overline{\underline{\mathrm{XII}}}$, p. 1 \ldots) pour conclure qu'un \unsure{homotopisme} est $\in W$, donc une catégorie contractile est $W$-asphérique, à fortiori une catégorie avec objet final (ou \unsure{avec|un} objet initial) est $W$-asphérique, \deleted{\unsure{D'où}} i.e. $\operatorname{Loc}(2)$.

\inserted{Comme $\operatorname{Loc}(1,2,3)\Longleftrightarrow\operatorname{Loc}(1,2\deleted{\unsure{a}},3\unsure{a|{}})$, il suffit de prouver \unsure{pour} $\operatorname{Loc}(2a)$.}

Prouvons $\operatorname{Loc}(3)$, \deleted{\unsure{donc il}} $\operatorname{Loc}(3\unsure{a|b})$ \unsure{va jouer|ce qui} \inserted{\unsure{a déjà}}. \deleted{Il suffit pour $X_{/S}$ \unsure{un objet}.} Soit $f:X\longrightarrow S$, $W$-\uline{co}asphérique, prouvons $f\in W$. Introduisons la catégorie \deleted{\unsure{co}}\uline{fibrée} sur $S$ dont les fibres sont les $ {}_{s\backslash}X_{\deleted{\unsure{s}}}$, soit $\Phi_S(X)$, et considérons l'inclusion canonique

\[
u_X:X\longrightarrow\Phi_S(X).
\]

\deleted{\uline{Lemme} Pour $X$ quelconque dans $\mathrm{Cat}_{/S}$, et moyennant $\operatorname{Loc}(1,2a,3')$, $i_X$ est}

% archive sheet number: 275

\scanpage{277}
\mspage{7}

\deleted{\unsure{On}} \deleted{Pour voir qu'}\inserted{Par $\operatorname{Loc}(3'\,b)$} elle est dans $W$.

\deleted{\unsure{Considérons|Remarquons}}

\deleted{\[
\unsure{i|\iota}_X^{o}:X^{o}\longrightarrow\Phi_{S^{o}}(X^{o}),
\]}

\deleted{elle est $\in W$ par $\operatorname{Loc}(3'\,b)$, donc si on prouve}

\deleted{\[
W=W^{o},\tag*{$(\unsure{*})$}
\]}

\deleted{on en conclura $i_X\in W$. \deleted{\unsure{que}} Admettons pour le moment que $\operatorname{Loc}(1,2a,3')$ implique $(*)$, et terminons de prouver qu'\unsure{elles|ils} impliquent $\operatorname{Loc}(3)$. \deleted{$\unsure{\Psi_S}$} Appliquons $\operatorname{Loc}(3')$ à $(\Psi_S(X^{\deleted{\unsure{o}}}))^o\longrightarrow S^{o}$, i.e. à $\Phi_{S^{o}}(X^{o})\longrightarrow S^{o}$, dont les fibres sont les $X_s^{o}$. Comme les $X_s$ sont $W$-asphériques, il résulte de \unsure{$(*)$|ceci|car} qu'il en est de $\unsure{=}$ des $X_s^{o}$, donc $\operatorname{Loc}(3')$ implique que $\Phi_{S^{o}}(X^{o})\longrightarrow S^{\unsure{o}}$ est dans $W$, donc \unsure{utilisant} \deleted{à \unsure{nouveau} \unsure{$(*)$|car} on conclut que}}

d'autre part, les fibres \inserted{$ {}_{s\backslash}X$} de $\Phi_S(X)$ sont $W$-coasphériques par hypothèse, donc par $\operatorname{Loc}(3'\,a)$ $\Phi_S(X)\unsure{{-}|\longrightarrow}S$ est dans $W$. Comme $i_X\in W$, il s'ensuit $(f:X{-}S)\in W$, qed.

\[
\text{b bis)}\qquad\operatorname{Loc}(1,2a,3'\,\mathrm{bis})\Longrightarrow\operatorname{Loc}(1,2,3\,\mathrm{bis}).
\]

On \deleted{sait} a déjà prouvé $\operatorname{Loc}(1,2,3a)$ \deleted{$\unsure{(3)}$} (vu que $\operatorname{Loc}(3'\,\mathrm{bis})$ contient $\operatorname{Loc}(3')$ comme cas particulier), \deleted{prouvons} donc $\operatorname{Loc}(3\,\mathrm{bis})$ en admettant déjà $\unsure{\Longrightarrow|{-}}\,3a$. Il suffira même de le prouver sous forme duale, puisque

\[
\operatorname{Loc}(1,2,3\,\mathrm{bis})\Longleftrightarrow\operatorname{Loc}(1,2^{o},3\,\mathrm{bis}^{o}),
\]

% archive sheet number: 276

\scanpage{278}
\mspage{8}

et que $\operatorname{Loc}(2^o)$ est connu par $\operatorname{Loc}(1,2,3)\Longleftrightarrow\operatorname{Loc}(1,2^o,3^o)$. Il reste donc à prouver que si

\[
f:X\unsure{{-}|\longrightarrow}Y\qquad\text{dans }\mathrm{Cat}_{/S}
\]

est telle que \deleted{$\unsure{{}_{s}f}$}

\[
{}_{s\backslash}f:{}_{s\backslash}X{-}{}_{s\backslash}Y
\]

est dans $W$ pour tt $s\in\operatorname{Ob}S$, alors $f\in W$. Considérons \unsure{alors|donc} le carré commutatif

\diagram{\begin{tikzcd}[ampersand replacement=\&]
X \arrow[r,"{\unsure{i_X|\iota_X}}"] \arrow[d,leftrightarrow,"f"'] \& \Phi_S(X) \arrow[d,"{\Phi_S(f)}"]\\
Y \arrow[r,"i_Y"] \& \Phi_S(Y),
\end{tikzcd}}

où $i_X,i_Y\in W$. par $\operatorname{Loc}(3'\,\mathrm{bis}\ b)$. L'hypothèse sur $f$ signifie que $\Phi_S(f)$ induit des $W$-isomorphismes sur les fibres, donc par $\operatorname{Loc}(3'\,\mathrm{bis})$ $\Phi_S(f)\in W$, donc $f\in W$ par $\operatorname{Loc}(1)$, qed.

\uline{2. Axiomes} \uline{Loc} (4) \uline{et} \uline{Loc} (5) (des sommes \unsure{directes|discrètes|disjointes}, de \unsure{connexité?})

% slanted
\margin{Il s'agit ici d'axiomes qu'on peut considérer comme étant de nature plus ou moins ``\unsure{triviale}''. Dans \unsure{tous les} cas à ma connaissance, ils sont bel et bien trivialement vérifiés.}

\uline{Loc 4} \inserted{(Axiome des sommes \unsure{discrètes|directes})} Pour une famille $(f_i)_{i\in I}$ de flèches de Cat, \unsure{si} $f=\coprod f_i$, on a $f\in W$ sss $f_i\in W$ pour tt $i\in I$.

Cet axiome se décompose en deux, en le ``\unsure{Si|si}'' et le ``seulement si'', que je vais \deleted{\unsure{désigner}} formuler sous la forme la plus économique :

\uline{Proposition} 2.\unsure{1|4}. L'axiome \uline{Loc} 4 \inserted{des sommes} \unsure{revient|équivaut} à la conjonction des deux axiomes suivants :

% archive sheet number: 277

\scanpage{279}
\mspage{9}

\uline{Loc} $4a$ : Soit $(f_i)_{i\in I}$ une famille de flèches de Cat, si les $f_i\in W$, alors $f\unsure{\in|=}\coprod f_i\in W$

\uline{Loc} $4b$ : Soit $S$ dans Cat, et $S_0\subset S$ un \unsure{sommand|sommant} direct (sous-catégorie à la fois ouverte et fermée). Soit $X$ dans $\mathrm{Cat}_{/S}$. Si $X\longrightarrow S$ est dans $W$, alors $X_0:=X\vert S_0\longrightarrow S_0$ aussi. (En d'autres termes, l'inclusion $S_0\unsure{\subset|\longrightarrow} S$ est une ``$W$-fibration'' i.e. \inserted{est} $\in\mathrm{Fib}_W$)

\uline{Corollaire} 2,2. Considérons la variante de \uline{Loc} $4a$ obtenue en se bornant aux familles finies. (Elle \unsure{s'étend|équivaut|s'ordonne} à \uline{Loc} 4 lorsque $W$ est stable par \uline{lim} filtrantes, cf axiome $\operatorname{Loc}(8)$ plus bas). Cette variante équivaut à

\uline{Loc} $4a_0$ : \deleted{\unsure{Soit}}\unsure{Tout|Soit} objet $X$ de Cat est ``$W$-cofibrant'' i.e. $\varnothing\longrightarrow X$ est dans $\mathrm{Cof}_W$, i.e. pour toute $f:Y\unsure{{-}|\longrightarrow}Y'$ qui est dans $W$, le morphisme correspondant $\mathrm{id}_X\amalg f:X\amalg Y{-}X\amalg Y'$ est dans $\underset{\inserted{\unsure{C}}}{\mathrm{Cof}}_W$.

\uline{Rappel} 2.3. De toutes façons, \deleted{pour tout}\inserted{si $W$ est un} \unsure{localiseur|locliseur} fondamental (ce que nous supposerons désormais, sauf mention du contraire), l'analogue \uline{fini} de \uline{Loc} $4a$ pour les produits est vrai, i.e. si \deleted{\unsure{les} $f$} on a une famille \uline{finie} de flèches $(f_i)_{i\in I}$ dans Cat, si les $f_i$ sont dans $W$, alors $f=\prod_i f_i\in W$. L'analogue pour une famille infinie n'est \inserted{vérifié} pratiquement \unsure{jamais pour un|pour un|jamais en un} localiseur fondamental.

% archive sheet number: 278

\scanpage{280}
\mspage{10}

\uline{Proposition} 2.\unsure{4|7}. Soit $W$ un localiseur fondamental. Alors les produits finis existent dans $\mathrm{Hot}_W$, et le foncteur $\mathrm{hot}_W:\mathrm{Cat}\unsure{{-}|\longrightarrow}\mathrm{Hot}_W\overset{\mathrm{def}}{=}(\mathrm{Cat})W^{-1}$ y commute. Si $W$ satisfait à \uline{Loc} \deleted{\unsure{3}}$4a$ (resp. \uline{Loc} $4a_0$), alors les sommes quelconques \inserted{(resp. finies)} existent dans $\mathrm{Hot}_W$, et le foncteur $\mathrm{hot}_W$ y commute. De plus, dans ce cas les sommes sont distributives p.r. aux produits (vu qu'il en est ainsi dans Cat).

La démonstration résulte simplement du résultat général :

\uline{Lemme} 2.5. Soit $(M,W)$ une catégorie de modèles ($W$ partie quelconque de $\mathrm{Fl}\,M$). Supposons \inserted{\unsure{que}} \deleted{toute famille \unsure{d'éléments}}\inserted{les produits \deleted{finis}} (resp. les produits \deleted{quelconques}\inserted{finis}) \inserted{(resp. les produits binaires)} existent dans $M$, et que de plus pour \deleted{\unsure{la|toute}} famille $(f_i)_{i\in I}$ de flèches de $M$, \deleted{telle} (resp tte famille finie, resp. \deleted{pour deux flèches}\inserted{famille binaire}), \deleted{\unsure{tels}} $f_i\in W$ pour tt $i\in I$ implique $f=\prod f_i\in W$,

Alors les produits \unsure{au type|du type|au degré} envisagé existent dans $\mathrm{Ho}_W=MW^{-1}$, et le foncteur \deleted{$\unsure{\mathrm{loc}_W}$} $\unsure{r_W|i_W}:M{-}MW^{-1}$ y commute. Enfin, si $M$ admet un objet final $e$, alors $e$ est aussi objet final dans $\mathrm{Ho}_W$.

% archive sheet number: 279

\scanpage{281}
\mspage{11}

On en \unsure{conclura|conclut} les énoncés duaux concernant les sommes directes et les objets initiaux.

La démonstration est un cas particulier du résultat bien utile suivant, appliqué au couple de foncteurs adjoints

\diagram{\begin{tikzcd}[ampersand replacement=\&,column sep=large]
M \arrow[r,shift left,"{\mathrm{diag}}"] \& M^I \arrow[l,shift left,"{\substack{\text{produits}\\\text{\unsure{dans|de} }M\\\text{indexés par }I}}" ] \& \mathpunct{:}
\end{tikzcd}}

% slanted
\margin{cf $\overline{\underline{\mathrm{XII}}}$ lemme \unsure{page} 248}

\uline{Lemme} 2.6. Soient $(M,W),(M',W')$ deux catégories de modèles, et

\diagram{\begin{tikzcd}[ampersand replacement=\&]
M \arrow[r,shift left,dash,"{\unsure{\varphi|\phi}}"] \& M' \arrow[l,shift left,"\psi"]
\end{tikzcd}}

un couple de foncteurs adjoints, compatibles aux localiseurs i.e. tels que

\[
\varphi(W)\subset W',\qquad \psi(W')\subset W.
\]

Alors les foncteurs \deleted{\unsure{induits}}\unsure{localisés}

\diagram{\begin{tikzcd}[ampersand replacement=\&]
MW^{-1} \arrow[r,shift left,"{\bar{\varphi}}"] \& M'W'^{-1} \arrow[l,shift left,"{\bar{\psi}}"]
\end{tikzcd}}

sont également adjoints, les flèches d'adjonction \inserted{pour $\bar{\varphi}\bar{\psi}$ et $\bar{\psi}\bar{\varphi}$} étant déduites de celles pour $\varphi\psi$ et $\psi\varphi$ par composition avec les foncteurs de localisation $M\unsure{{-}|\longrightarrow}MW^{-1}$, $M'{-}M'W'^{-1}$.

% archive sheet number: 280

\scanpage{282}
\mspage{12}

\uline{Loc} (5) \inserted{(Axiome de \uline{connexité}) ou de \uline{non-trivialité})} Soit $\unsure{f}\in W$, alors $\pi_0(f)$ est bijectif.

Cette \deleted{axiome} propriété de $W$ est d'une \unsure{utilité} \unsure{manifeste|constante} en pratique, quand on travaille avec $W$ au lieu de se borner à construire des $W$ et faire joujou avec les axiomes et leur stabilité par passage à des $\underline{\operatorname{Hom}}(I,M)$ (où ici $M=\mathrm{Cat}$). D'autre part, on peut aussi la regarder comme un axiome de \uline{non-trivialité} pour $\mathrm{Hot}_W$. Je rappelle à ce sujet : \deleted{\illegible{1}}

\uline{Proposition} 2.7. Soit $W$ un localiseur fondamental. Conditions équivalentes

a) $W$ satisfait l'axiome de connexité \uline{Loc} 5

b) Le foncteur $\pi_0:\mathrm{Cat}\unsure{{-}|\longrightarrow}\mathrm{Ens}$ se factorise par un foncteur

\diagram{\begin{tikzcd}[ampersand replacement=\&,column sep=large]
\mathrm{Cat} \arrow[r,dotted] \arrow[rr,bend right=24,dashed,"{\pi_0}"'] \& \mathrm{Hot}_W \arrow[r,"{\pi_0}"] \& \mathrm{Ens},
\end{tikzcd}}

% slanted
\margin{\uline{NB} Si $W$ est fort saturé (condition $\operatorname{Loc}(1\,\mathrm{ter})$), alors les conditions ci-contre (i.e. $\operatorname{Loc}(5)$) équivalent à ceci : \deleted{\unsure{b|a}}\unsure{d)} $W$ n'est égal ni à $\mathrm{Fl}(\mathrm{Cat})$ tout entier, ni à $\mathrm{Fl}(\mathrm{Cat}^{*})\cup\{\mathrm{id}_{\varnothing}\}$ (i.e. $\mathrm{Fl}(\mathrm{Cat})\setminus W$ n'est ni vide, ni égal à l'ensemble des flèches $\varnothing\longrightarrow X$ où $X\ne\varnothing$). Si $W$ satisfait à $\operatorname{Loc}(\unsure{7\,D_{W_{\omega}}|7\,D_{W\omega}|7\,D_{Ww}})$, alors $W$ \deleted{ne} satisfait $\operatorname{Loc}(5)$ sss $W\ne\mathrm{Fl}(\mathrm{Cat})$ (cf XII th. 2, p. 20)}

c) La catégorie $\mathrm{Hot}_W^{*}=\mathrm{Hot}_W\setminus\{\varnothing\}$ n'est équivalente à la catégorie ponctuelle

La catégorie $\mathrm{Hot}_W$ n'est pas équivalente \deleted{\unsure{à}} ni à $\unsure{\mathbb{\Delta}|\Delta}^{1}$.

Pour la démonstration et pour d'autres formes équivalentes de l'axiome de \unsure{connexité},

% archive sheet number: 281

\scanpage{283}
\mspage{13}

on en déduirait $\pi_0(X)\longrightarrow\varnothing$, \unsure{ou|or|i.e.} $\pi_0(X)=\varnothing$, absurde !.

Reste le cas où $X$ est non vide. Cela résulte du \deleted{fait \unsure{général} \unsure{telle} $X$ \unsure{ayant} une section}\inserted{fait général}

\uline{Lemme} \unsure{2|?}.\unsure{13|14}. Soit \deleted{$\unsure{C}$}$M$ une catégorie ayant un objet initial \uline{strict} $\varnothing_M$ (\deleted{si \unsure{le}} $\operatorname{Hom}_M(X,\varnothing_M)\ne\varnothing$, on a $X=\varnothing_M$), \inserted{un objet final $e_M$,} et telle que pour tt $X\ne\varnothing_M$, $X\longrightarrow e$ ait une section i.e. $\operatorname{Hom}(e,X)\ne\varnothing$.

\deleted{Alors} Soit $W\subset\mathrm{Fl}\,C$, \deleted{Alors tel que \unsure{pour}} \inserted{\textcircled{1\textsuperscript{o}}} ne contenant pas de flèches $\varnothing_M\longrightarrow X$, avec $X\ne\varnothing_M$, et \inserted{\textcircled{2\textsuperscript{o}}} tel que

% margin note written sideways
\inserted{pour tout diagramme
\diagram{\begin{tikzcd}[ampersand replacement=\&]
X_1 \& X_2 \arrow[l,"u_2"']\\
e \arrow[u,"u_1"] \arrow[ur,dashed,"s"'] \&
\end{tikzcd}}
\unsure{avec|où} $u_2\in W$ \deleted{\diagram{\begin{tikzcd}[ampersand replacement=\&]
\unsure{X_1} \& \unsure{X_2} \arrow[l,"{\unsure{u_2}}"']\\
\unsure{e} \arrow[u,"{\unsure{u_1}}"] \arrow[ur,dotted,"{\unsure{s}}"'] \&
\end{tikzcd}} \unsure{avec}} $\exists\,s$ \uline{dans $M$} \inserted{tel} que $u_2s\unsure{=|\simeq}u_1$ dans $MW^{-1}$.}

a) $\varnothing$ est un objet initial \uline{strict} de $MW^{-1}$

b) Pour tt $X\ne\varnothing$, $\operatorname{Hom}_M(\deleted{\unsure{X}}e,\deleted{\unsure{e}}X)\unsure{{-}|\longrightarrow}\operatorname{Hom}_W(e,X)$ est surjectif.

\uline{Dém} Que $\varnothing_M$ soit objet initial de $MW^{-1}$ a été vu (dernière assertion du lemme 2.5, forme duale), prouvons qu'il est strict, i.e. que si \deleted{$\unsure{X\to\varnothing_M}$ \unsure{dans}} $\operatorname{Hom}_W(X,\varnothing)\ne\varnothing$, \inserted{on a $X=\varnothing_M$.} En effet, \deleted{\unsure{sinon}} si on avait $u:X{-}\varnothing$ dans $MW^{-1}$, $u$ \deleted{\unsure{ser}}serait représenté par

\diagram{\begin{tikzcd}[ampersand replacement=\&,column sep=small]
X=X_0 \& X_1 \arrow[l,"W"'] \arrow[r] \& X_2 \& \cdots \arrow[l,"W"'] \& X_{2n-2}\arrow[r] \& X_{2\unsure{n}-1} \& X_{2n}=\varnothing_M \arrow[l,"W"']
\end{tikzcd}}

On procède par récurrence sur la \inserted{demi-}longueur de cette chaîne pour voir que $X=\varnothing$. En effet, l'hyp. sur $W$ implique que $X_{2n-1}=\varnothing_M$, puis celle que $\varnothing_M$ est strict implique $X_{\unsure{2-n|2n-2}}=\varnothing_M$, d'où la conclusion.

% archive sheet number: 282

\scanpage{284}
\mspage{14}

\deleted{Soit} Prouvons la surjectivité de $\operatorname{Hom}_M(e,X)\unsure{{-}|\longrightarrow}\operatorname{Hom}_W(e,X)$ pour $X\ne\varnothing_M$, en montrant que tout chemin réduit de longueur $2n+1$ est ``homotope'' à un chemin réduit de longueur $2n-1$ :

\diagram{\begin{tikzcd}[ampersand replacement=\&,column sep=small]
e=X_0 \arrow[r,"{\deleted{\unsure{\leftarrow}}}" description] \& X_{\unsure{1|p}} \& X_2 \arrow[l,"W"'," {\deleted{\unsure{\rightarrow}}}" description] \arrow[r,dash] \& X_3 \& \cdots \arrow[l] \& X_{2n} \arrow[l,"{\deleted{\unsure{\rightarrow}}}" description] \arrow[r,"{\deleted{\unsure{\leftarrow}}}" description] \& X_{2n+1}=X
\end{tikzcd}}

En effet, on peut remplacer la section

\diagram{\begin{tikzcd}[ampersand replacement=\&,column sep=small]
e=X_0 \arrow[r,"{\deleted{\unsure{W}}}","u_1"'," {\deleted{\unsure{\leftarrow}}}" description] \arrow[rr,bend right=40,dashed,"s"'] \arrow[r,bend right=22,"{\deleted{\unsure{s}}}"'] \& X_1 \& X_2 \arrow[l,"W"',"u_2"] \arrow[r,"{\deleted{\unsure{W}}}","u_3"'] \& X_3 \& X_4 \arrow[l,"W"']
\end{tikzcd}}

\deleted{\unsure{puis} d'abord par}

\deleted{\diagram{\begin{tikzcd}[ampersand replacement=\&]
\unsure{e} \arrow[r,"{\unsure{u_2s}}"] \& \unsure{X_2} \& \unsure{X_3}\arrow[l,"{\unsure{W}}"'] \& \cdots
\end{tikzcd}}}

par

\diagram{\begin{tikzcd}[ampersand replacement=\&]
e\arrow[r,"{\deleted{\unsure{u_3s}}}","{u_3s}"'] \& X_3 \& \cdots \arrow[l,"W"']
\end{tikzcd}}

\deleted{(où $s$ est une section de $X_1$ sur $e=X_0$), puis par}

où $s$ est une section de $X_2$ sur $e$ telle que

\[
u_2s\deleted{\unsure{\simeq}}=u_1\qquad\text{\uuline{dans }}MW^{-1}.
\]

Cela prouve donc par \unsure{réc} sur $n$ que le chemin est homotope à un chemin de longueur $1$, $e\longrightarrow X$, qed.

Il faut donc vérifier, \deleted{1\textsuperscript{o}} dans le cas actuel, vérifier \uline{1\textsuperscript{o}}) (c'est le \inserted{cor.} \unsure{cor.} \unsure{2|?}.12.), et 2\textsuperscript{o})

\diagram{\begin{tikzcd}[ampersand replacement=\&]
X_1 \& X_2 \arrow[l,"{u_2\in W}"']\\
e \arrow[u,"u_1"] \arrow[ur,dotted,"s"'] \&
\end{tikzcd}}

Mais comme $u_2\in W$, on a $\pi_0(u_2)$ \unsure{iso|bij.}, donc

% archive sheet number: 283

\scanpage{285}
\mspage{13}

je renvoie à $\overline{\underline{\unsure{\mathrm{XIV}|\mathrm{XV}}}}$ \deleted{\unsure{8|3}}9.12 \inserted{(p. 133 bis).} Pour le nom ``axiome de connexité'', je peux offrir l'explication suivante 2.8. et 2.9.

\uline{Corollaire} 2.8. Si \deleted{\unsure{le} \unsure{est}} l'axiome de connexité $\operatorname{Loc}(5)$ est satisfait, alors les objets $W$-asphériques sont $0$-connexes (tautologique), et la réciproque est \unsure{vraie} si $W$ est fortement saturé, \deleted{ou si $\operatorname{Loc}(4b)$ est satisfait.}

\inserted{(ou aussi, si $W$ satisfait l'axiome des fibrations $\operatorname{Loc}\,\unsure{7}$ ci-dessous, puisque par $\overline{\underline{\mathrm{XII}}}$ th. 2 p. \unsure{21|31}, dans ce cas, $\operatorname{Loc}(5)$ équivaut à $W\ne\mathrm{Fl}(\unsure{M|\mathrm{Cat}})$ \unsure{\ldots})}

Car si l'axiome de connexité n'est pas satisfait, alors pour tt $X$ non vide \unsure{en|de} Cat, $X\unsure{{-}|\longrightarrow}e$ est iso, donc ($W$ étant fort\textsuperscript{mt} saturé) $X{-}e$ est dans $W$ i.e. $X$ est \deleted{dans} $W$-asphérique, donc il y a des objets \unsure{$W$|hot}{}-asphériques non $0$-connexes.

\deleted{Dans le cas où $\operatorname{Loc}(4b)$ est satisfait, et que tout objet $W$-asphé}

Notons aussi

\uline{Corollaire} 2.\unsure{9|8}. Supposons que $W$ satisfait à l'axiome $\operatorname{Loc}(4)$ des sommes discrètes (ou seulement à $\operatorname{Loc}\,4b)$. Pour que $W$ satisfasse l'axiome de connexité il suffit que pour tte flèche $f:X{-}Y$ avec $Y$ $0$-connexe, si $f$ est \deleted{\unsure{de}} $\in W$ alors $X$ est $0$-connexe. (C'est nécessaire bien sûr \unsure{sans|dans} \inserted{supposer} $\operatorname{Loc}(4b)$ ni rien \unsure{du|de} \deleted{\unsure{tout}}\unsure{tel|tout})

% archive sheet number: 284

\scanpage{286}
\mspage{14}

2.7. \deleted{\unsure{8|7}.} Considérons maintenant le foncteur composé

\diagram{\[
\begin{tikzcd}[ampersand replacement=\&]
\mathrm{Ens} \arrow[r,"{I\mapsto I_{\mathrm{Cat}}}"] \& \mathrm{Cat} \arrow[r,"{\mathrm{hot}_W}"] \& \mathrm{Hot}_W,
\end{tikzcd}
\qquad I\longrightarrow \substack{I_{\mathrm{Hot}_W}\\\text{ou } I_W},
\tag*{(2.7.1)}
\]}

où le premier foncteur est celui qui associe à l'ens $I$ la catégorie discrète $I_{\mathrm{Cat}}$ ayant $I$ comme ens d'objets. (Le plus souvent, on identifie $I_{\mathrm{Cat}}$ avec $I$ dans les notations, tout comme pour un ens ordonné $J$, on identifie $J$ et la ``catégorie ordonnée'' définie par $J$. Ici $I$ peut être considéré comme muni de l'ordre discret.) On a d'ailleurs

\[
 I_{\mathrm{Cat}}\simeq I\times e\quad (\deleted{\unsure{\mathrm{i.e.}}}\overset{\mathrm{def}}{=}\text{ somme de }I\text{ copies de }e),
 \tag*{(2.7.2)}
\]

donc quand \deleted{$\mathrm{hot}_W$} $\operatorname{Loc}(4a)$ est satisfait \deleted{\unsure{alors}}\unsure{donc|alors} $\mathrm{hot}_W$ commute aux sommes \unsure{directes|discrètes} on aura un isom canonique

\[
 I_W\simeq I\times e_W. \tag*{(2.7.3.)}
\]

(où $e_W$ est l'objet final de $\mathrm{Hot}_W$, i.e. $e$ considéré comme objet de $\mathrm{Hot}_W$).

\uline{Prop}. 2.8. Supposons \deleted{que} l'axiome \deleted{$\operatorname{Loc}\,\unsure{3|5}$} de connexité $\operatorname{Loc}(5)$ satisfait, d'où un foncteur

\[
 \pi_0:\mathrm{Hot}_W\longrightarrow\mathrm{Ens}. \tag*{(2.8.1)}
\]

Ce foncteur est \uline{adjoint à gauche} du foncteur $I\mapsto I_W$ (2.7.1), i.e. on a une bijection fonctorielle en $(X,I)$ de $\mathrm{Hot}_W\times\mathrm{Ens}$ :

% archive sheet number: 285

\scanpage{287}
\mspage{15}

\[
 \operatorname{Hom}_{\mathrm{Ens}}(\pi_0(X),I)\simeq\operatorname{Hom}_W(X,I_W).
 \tag*{(2.8.2)}
\]

\deleted{Quand $W$ satisfait l'axiome $\operatorname{Loc}(4a)$ des sommes \unsure{d'ordre|directes}, alors} \inserted{Le foncteur} $\pi_0$ est également \uline{adjoint à droite} du \uline{même} foncteur $I\mapsto I_W$, i.e. on a une bijection fonctorielle

\[
 \operatorname{Hom}_{\mathrm{Ens}}(I,\pi_0(X))\simeq\operatorname{Hom}_W(I_W,X)\unsure{;|.}
 \tag*{(2.8.3)}
\]

Prenant $I=e$ (ensemble ponctuel), on trouve le \deleted{corollaire :}

\uline{Corollaire} 2.9. \deleted{Moyennant $\operatorname{Loc}(4a)$,} On a \deleted{\unsure{la bijec|la bijection}} l'isom fonctoriel en $X$ de $\mathrm{Hot}_W$

\[
 \boxed{\operatorname{Hom}_W(e_W,X)\simeq\pi_0(X)}. \tag*{(2.8.4)}
\]

\uline{Commentaire} 2.10. La formule d'adjonction (2.8.2) est d'ailleurs familière, chaque fois qu'on a une catégorie $\unsure{H|\mathcal{H}}$ (telle $\mathrm{Cat}$, ou $\mathrm{Hot}_W$ etc) munie d'un foncteur $\pi_0$, \inserted{$\pi_0(X)$} correspondant à l'intuition d'une décomposition d'un objet en ``composantes connexes''. (Je ne vais pas expliciter ici les conditions \inserted{\unsure{générales|génériques|formelles}} dans lesquelles on \unsure{définit|introduit|entend} un tel foncteur.) Par contre, l'adjonction en sens inverse de ces deux mêmes \unsure{foncteurs} est un phénomène tout à fait spécial à des catégories

% archive sheet number: 286
